% chapters/02_telemetry.tex

\chapter{Falha de Telemetria}
\epigraph{``Há uma rachadura em tudo.\\É assim que a luz entra.''}{Leonard Cohen, "Anthem"}

O transurbano parou com um chiado hidráulico na borda do Distrito Industrial, um som que lembrava um suspiro de exaustão. Elara desembarcou em uma paisagem de aço corrugado enferrujado e asfalto rachado. Aqui, o relevo de São Paulo se achatava em uma planície recuperada, um lugar onde a ambição brilhante da cidade havia estagnado décadas atrás.

Ela estava parada na esquina da Radial Leste com a Salim Farah Maluf, sua respiração vindo em goles rasos e cautelosos. \textit{Não pense no inalador,} disse a si mesma. \textit{Pensar no ar faz você precisar mais dele.} 

Era um truque que aprendera aos vinte e dois anos: trate seu corpo como uma máquina falhando. Se você não reconhecer o atrito, talvez as engrenagens continuem girando por mais um turno.

\systemcall{LOCALIZAÇÃO: SETOR 7-B // STATUS: UTILIDADE SUB-OTIMIZADA \\ AVISO: DEGRADAÇÃO DE SINAL DETECTADA // AUDITORIA NECESSÁRIA}

A lente Aegis-7 oscilou, um ícone vermelho nervoso pulsando no canto de seu olho. O HUD corporativo odiava este lugar. Para os algoritmos do Auditor, o Setor 7-B era um incômodo estatístico --- uma coleção de ``ativos zumbis'' que se recusavam a ser produtivos ou a morrer completamente.

Ela começou a caminhar. O silêncio aqui era diferente. No Centro, o silêncio era caro, um produto gerenciado de vidros acústicos e filtros de alta tecnologia. Aqui, o silêncio era pesado, um vácuo deixado por um mercado que havia simplesmente ido embora.

\textit{Este é o lugar,} pensou ela, parando diante do galpão cinzento que vira da janela do ônibus. A porta de metal estava marcada com o logo desbotado de uma logística de entrega que não existia mais. No HUD, o prédio brilhava com um aviso roxo.

\systemcall{ERRO DE PROTOCOLO: DISCREPÂNCIA DE DADOS DETECTADA \\ O SENSOR DE PRESENÇA 092 RELATA: VAZIO \\ O FLUXO DE CALOR RELATA: ATIVIDADE BIOLÓGICA}

Elara hesitou. O protocolo era claro: documentar a falha do sensor, recalibrar o nó e retornar à Paulista para o processamento de fim de turno. Mas o som --- aquela ressonância de baixa frequência que ela sentira no ônibus --- era mais forte aqui. Vinha de trás da porta.

Ela estendeu a mão, e por um momento, a lente Aegis-7 gritou um aviso de ``Violação de Perímetro'' com uma voz que soava como a de sua mãe. \textit{Isso é uma brecha. Se eu ficar, o Auditor verá a desconexão. Eles vão sinalizar meu nível de saúde. Eles vão—}

A porta rangeu. Não muito --- apenas uma fresta.

Um homem estava lá. Ele não vestia um terno corporativo e não tinha os olhos assombrados e vazios dos ``Não-Assinantes'' que dormiam nos túneis do metrô. Ele parecia saudável. Obscenamente saudável. Sua pele tinha um brilho que não deveria existir em uma cidade que não via o sol há três semanas.

Ele não segurava uma arma. Segurava uma pequena tigela de madeira esculpida à mão.

``Você está atrasada, Elara'', disse ele. Sua voz era calma, desprovida daquela urgência transacional que ela ouvia em todos os outros. 

\textit{Ele sabe meu nome.} O pensamento deveria tê-la aterrorizado, mas, em vez disso, pareceu uma chave girando em uma fechadura enferrujada. \textit{Como ele pode saber meu nome? Sou apenas uma variável. Sou apenas uma Analista Júnior.}

``Eu... eu estou aqui para a auditoria dos sensores'', ela gaguejou, sua mão indo instintivamente à Aegis-7, tentando forçá-la a voltar online. 

O homem sorriu, uma expressão lenta e triste. ``Os sensores não estão quebrados, Elara. Eles só estão procurando por coisas que não estão mais aqui. Você gostaria de ver o que realmente está aqui?''

Ela olhou para trás, para a Garoa cinzenta e sufocante da rua, e depois para o calor âmbar que escapava pela fresta da porta. Pela primeira vez, a escolha não pareceu um cálculo de créditos.
